\documentclass[11pt]{article}
\usepackage[margin=1in]{geometry}
\usepackage{amsmath,amssymb,amsthm}

\newtheorem{theorem}{Theorem}
\newtheorem{proposition}[theorem]{Proposition}
\newtheorem{lemma}[theorem]{Lemma}
\newtheorem{corollary}[theorem]{Corollary}
\theoremstyle{definition}
\newtheorem{remark}[theorem]{Remark}

\newcommand{\R}{\mathbb{R}}
\newcommand{\abs}[1]{\left|#1\right|}
\newcommand{\dx}{\,dx}
\DeclareMathOperator{\Asym}{\mathcal{A}}

\title{The Kohler--Jobin transfer:\\
sharp Faber--Krahn stability from sharp Saint--Venant stability}
\author{Plan 1 / Route A, building block \textsc{KohlerJobinTransfer}}
\date{}

\begin{document}
\maketitle

\begin{abstract}
We record the Kohler--Jobin transfer step.
\emph{Input:} the sharp Saint--Venant stability inequality
\[
E(\Omega)-E(B^*)\ge c_{\mathrm{SV}}(N)\,
\Bigl(\frac{|B^*|}{|B_1|}\Bigr)^{(N+2)/N}\Asym(\Omega)^{2}
\qquad\text{for every open }\Omega,\ \abs{\Omega}=\abs{B^*},
\]
produced by the \textsc{BoundedReduction} block from the sibling
\textsc{SaintVenantAssembly} block.
\emph{Output:} the sharp Faber--Krahn stability inequality
\[
\Bigl(\frac{|\Omega|}{|B_1|}\Bigr)^{2/N}
\bigl(\lambda_1(\Omega)-\lambda_1(B^*)\bigr)
\ge c_{\mathrm{FK}}(N)\,\Asym(\Omega)^{2}.
\]
The argument is non-perturbative: it uses only the Kohler--Jobin
inequality $T\cdot\lambda_1^{(N+2)/2}\ge T_N L_N^{(N+2)/2}$ (cited from
Brasco~\cite{Brasco2014}) and one elementary one-variable inequality.
No compactness, no contradiction, no Taylor expansion.
\end{abstract}

\section{Notation}\label{sec:notation}

For open bounded $\Omega\subset\R^N$ ($N\ge2$):
\begin{itemize}
\item $u_\Omega\in H_0^1(\Omega)$ is the torsion potential, the unique
weak solution of $-\Delta u=1$ in $\Omega$, $u=0$ on $\partial\Omega$;
\item $T(\Omega)=\int_\Omega u_\Omega\dx$ is the torsional rigidity;
\item $\lambda_1(\Omega)$ is the first Dirichlet eigenvalue;
\item $E(\Omega):=-T(\Omega)/2$ is the BDV-normalized Saint--Venant
energy, so $T(B^*)-T(\Omega)=2(E(\Omega)-E(B^*))$;
\item $B^*$ is the ball with $|B^*|=|\Omega|$, $B_R$ is the ball of
radius $R$ at the origin;
\item $\Asym(\Omega):=\inf_{x_0\in\R^N}|\Omega\Delta(x_0+B^*)|/|B^*|$ is
the Fraenkel asymmetry;
\item $B_1$ is the unit ball, $L_N:=\lambda_1(B_1)$, $T_N:=T(B_1)=|B_1|/(N(N+2))$.
\end{itemize}

The scaling laws
\begin{equation}\label{eq:scaling}
T(t\Omega)=t^{N+2}T(\Omega),\qquad
\lambda_1(t\Omega)=t^{-2}\lambda_1(\Omega)\qquad(t>0),
\end{equation}
make $\Psi(\Omega):=T(\Omega)\cdot\lambda_1(\Omega)^{(N+2)/2}$ scale
invariant. Write $\beta:=(N+2)/2$, $L:=\lambda_1(B^*)$, $T_0:=T(B^*)$.

\section{The Kohler--Jobin inequality}\label{sec:kj}

\begin{theorem}[Kohler--Jobin, 1978]\label{thm:kj}
For every open bounded $\Omega\subset\R^N$ with $\lambda_1(\Omega)<\infty$
and every ball $B$,
\begin{equation}\label{eq:kjglobal}
T(\Omega)\cdot\lambda_1(\Omega)^{\beta}
\ge T(B)\cdot\lambda_1(B)^{\beta}.
\end{equation}
Equivalently $\Psi(\Omega)\ge\Psi(B_1)=T_N L_N^{\beta}$, with equality
iff $\Omega$ is a ball (up to translation and a null set).
\end{theorem}

We invoke~\eqref{eq:kjglobal} as a black box; for a modern proof by
Schwarz symmetrization see Brasco~\cite[Theorem~3.3]{Brasco2014}.

\begin{remark}[Scaling check]\label{rem:scaling}
$\beta=(N+2)/2$ is forced by scale-invariance of~\eqref{eq:kjglobal}.
From~\eqref{eq:scaling}, $T(t\Omega)\lambda_1(t\Omega)^\beta =
t^{N+2-2\beta}T(\Omega)\lambda_1(\Omega)^\beta$; invariance forces
$N+2-2\beta=0$.
\end{remark}

\section{Pointwise transfer at fixed volume}

\begin{lemma}[Pointwise Kohler--Jobin transfer]\label{lem:pointwise}
For every open $\Omega\subset\R^N$ with $|\Omega|=|B^*|$,
\begin{equation}\label{eq:ptwise}
\lambda_1(\Omega)\ge L\cdot\Bigl(\frac{T_0}{T(\Omega)}\Bigr)^{1/\beta}
=L\cdot\Bigl(\frac{T_0}{T(\Omega)}\Bigr)^{2/(N+2)}.
\end{equation}
\end{lemma}

\begin{proof}
Apply Theorem~\ref{thm:kj} with $B=B^*$ and rearrange.
\end{proof}

By the classical Saint--Venant inequality $T(\Omega)\le T_0$ at fixed
volume, the Saint--Venant deficit
$\delta_{\mathrm{SV}}(\Omega):=T_0-T(\Omega)\ge0$ is non-negative. Set
$s:=\delta_{\mathrm{SV}}/T_0\in[0,1)$; \eqref{eq:ptwise} becomes
\begin{equation}\label{eq:ptwises}
\lambda_1(\Omega)-L\ge L\bigl[(1-s)^{-1/\beta}-1\bigr].
\end{equation}

\section{An elementary one-variable inequality}

\begin{lemma}\label{lem:onevar}
For $p\in(0,1]$, $s\in[0,1)$,
\begin{equation}\label{eq:onevar}
(1-s)^{-p}-1\ge p\,s.
\end{equation}
\end{lemma}

\begin{proof}
$f(s):=(1-s)^{-p}-1$ satisfies $f(0)=0$ and
$f'(s)=p(1-s)^{-(p+1)}\ge p$ for $s\ge0$. By FTC,
$f(s)=\int_0^s f'(\tau)d\tau\ge ps$.
\end{proof}

\begin{remark}
No Taylor expansion; the inequality is exact on $[0,1)$. The factor $p$
is sharp at $s=0$ since $f'(0)=p$. We use $p=1/\beta=2/(N+2)\in(0,1]$.
\end{remark}

\section{Small-deficit linearization}

\begin{lemma}[Small-deficit linearization]\label{lem:linear}
For every open $\Omega\subset\R^N$ with $|\Omega|=|B^*|$ and
$\delta_{\mathrm{SV}}(\Omega)\le T_0/2$,
\begin{equation}\label{eq:linear}
\lambda_1(\Omega)-L\ge\frac{2L}{(N+2)T_0}\,\delta_{\mathrm{SV}}(\Omega).
\end{equation}
\end{lemma}

\begin{proof}
$s=\delta_{\mathrm{SV}}/T_0\le1/2$, $p=2/(N+2)$. By \eqref{eq:ptwises}
and Lemma~\ref{lem:onevar},
$\lambda_1(\Omega)-L\ge L[(1-s)^{-p}-1]\ge Lps=Lp\delta_{\mathrm{SV}}/T_0
=(2L/((N+2)T_0))\delta_{\mathrm{SV}}$.
\end{proof}

The constant $2L/((N+2)T_0)$ is the universal \emph{Kohler--Jobin
multiplier}. Both sides of~\eqref{eq:linear} scale, so it depends only
on $N$: $2L/((N+2)T_0)=2L_N/((N+2)T_N)=2NL_N/|B_1|$.

\section{Large-deficit branch}

\begin{lemma}[Large-deficit bound]\label{lem:large}
If $|\Omega|=|B^*|$ and $\delta_{\mathrm{SV}}(\Omega)>T_0/2$,
\begin{equation}\label{eq:large}
\lambda_1(\Omega)-L\ge L\bigl(2^{2/(N+2)}-1\bigr).
\end{equation}
\end{lemma}

\begin{proof}
$T_0/T(\Omega)>2$, so~\eqref{eq:ptwise} gives
$\lambda_1(\Omega)\ge L\cdot 2^{2/(N+2)}$.
\end{proof}

Since $\Asym(\Omega)\in[0,2)$ for $|\Omega|=|B^*|$,
$\Asym(\Omega)^2\le4$, so~\eqref{eq:large} yields in this regime
\begin{equation}\label{eq:largeA}
\lambda_1(\Omega)-L\ge\frac{L(2^{2/(N+2)}-1)}{4}\,\Asym(\Omega)^2.
\end{equation}

\section{Main transfer theorem}\label{sec:main}

\begin{theorem}[Sharp Faber--Krahn via Kohler--Jobin]\label{thm:main}
Assume sharp Saint--Venant stability in energy form: there exists
$c_{\mathrm{SV}}(N)>0$
such that
\begin{equation}\label{eq:svinput}
E(\Omega)-E(B^*)\ge c_{\mathrm{SV}}(N)
\Bigl(\frac{|B^*|}{|B_1|}\Bigr)^{(N+2)/N}\Asym(\Omega)^2
\quad\forall\,\Omega,\ 0<|\Omega|=|B^*|<\infty.
\end{equation}
Then for every such $\Omega$,
\begin{equation}\label{eq:fkout}
\Bigl(\frac{|\Omega|}{|B_1|}\Bigr)^{2/N}
\bigl(\lambda_1(\Omega)-\lambda_1(B^*)\bigr)
\ge c_{\mathrm{FK}}(N)\,\Asym(\Omega)^2,
\end{equation}
with
\begin{equation}\label{eq:cfk}
c_{\mathrm{FK}}(N):=\min\!\left\{
\frac{4L_N\,c_{\mathrm{SV}}(N)}{(N+2)T_N},\;\;
\frac{L_N(2^{2/(N+2)}-1)}{4}
\right\}.
\end{equation}
\end{theorem}

\begin{proof}
Split on $\delta_{\mathrm{SV}}\le T_0/2$ vs.\ $\delta_{\mathrm{SV}}>T_0/2$.

\emph{Case 1.} Put $r=(|B^*|/|B_1|)^{1/N}$. Lemma~\ref{lem:linear} and
\eqref{eq:svinput} give
$\lambda_1(\Omega)-L\ge(2L/((N+2)T_0))\delta_{\mathrm{SV}}
\ge(4L\,c_{\mathrm{SV}}r^{N+2}/((N+2)T_0))\Asym^2
=(4L_N c_{\mathrm{SV}}/((N+2)T_N))r^{-2}\Asym^2$ by scaling. Multiplying
by $r^2=(|\Omega|/|B_1|)^{2/N}$ gives the first constant in~\eqref{eq:cfk}.

\emph{Case 2.} \eqref{eq:largeA} gives
$\lambda_1(\Omega)-L\ge r^{-2}(L_N(2^{2/(N+2)}-1)/4)\Asym^2$ by scaling.
Multiplying by $r^2$ gives the second constant.

Either way, the conclusion holds with~\eqref{eq:cfk}.
\end{proof}

\begin{remark}[Energy normalization]
With $E:=-T/2$, the torsion deficit is
$\delta_{\mathrm{SV}}=T(B^*)-T(\Omega)=2(E(\Omega)-E(B^*))$. This is the
source of the factor $4L_N/((N+2)T_N)$ in the small-deficit branch of
\eqref{eq:cfk}.
\end{remark}

\begin{remark}[Sharpness]
The exponent~2 on $\Asym$ in~\eqref{eq:fkout} is sharp; it is saturated by
smooth ellipsoidal perturbations of the ball where
$\lambda_1(\Omega)-L$ and $\Asym^2$ are of the same order
(\cite{Brasco2014} and references).
\end{remark}

\section{Summary of constants}\label{sec:const}

\begin{center}
\renewcommand{\arraystretch}{1.25}
\begin{tabular}{|l|l|l|}
\hline
Symbol & Value / dependence & Source \\
\hline
$N$ & dimension, $\ge2$ & hypothesis \\
$L_N$ & $\lambda_1(B_1)=j_{N/2-1,1}^2$ & \\
$T_N$ & $T(B_1)=|B_1|/(N(N+2))$ & \\
$\beta=(N+2)/2$ & scaling exponent & Rmk.~\ref{rem:scaling} \\
$L,T_0$ & $\lambda_1(B^*),T(B^*)$ & \\
$2L_N/((N+2)T_N)=2NL_N/|B_1|$ & KJ multiplier & Lem.~\ref{lem:linear} \\
$L_N(2^{2/(N+2)}-1)$ & large-deficit gap & Lem.~\ref{lem:large} \\
$c_{\mathrm{SV}}(N)$ & global sharp Saint--Venant energy constant & \textsc{BoundedReduction} \\
$c_1(N)$ & $4L_N c_{\mathrm{SV}}/((N+2)T_N)$ & small branch \\
$c_2(N)$ & $L_N(2^{2/(N+2)}-1)/4$ & large branch \\
$c_{\mathrm{FK}}(N)$ & $\min\{c_1(N),c_2(N)\}$ & Thm.~\ref{thm:main} \\
\hline
\end{tabular}
\end{center}

All constants are explicit; none depends on a compactness or contradiction
step in the Kohler--Jobin transfer.

\begin{thebibliography}{9}
\bibitem{Brasco2014} L.~Brasco,
\emph{On torsional rigidity and principal frequencies: an invitation to
the Kohler--Jobin rearrangement technique},
ESAIM Control Optim. Calc. Var. \textbf{20} (2014), 315--338.
\end{thebibliography}

\end{document}
